% ============================================================
% 1.2 INTRODUCTION: THE FAIRNESS QUESTION -- Writer: Adam Bajer
% Target: about 300 words. Comes right after Ignacio's 1.1 (backgrounds and Figure 1); read it first so the two parts connect.
% Brief: "Introduction, including Background and Related Work: research
% question, motivation, dataset significance, objectives, and the relevant
% scientific or technical background and prior work ... integrated into the
% Introduction rather than submitted as a separate chapter."
% IRENE'S FEEDBACK ON THE SENTENCE OUTLINE (follow it in this order):
%   "Begin with how students' backgrounds can shape their choice of degree,
%    then explain why GPA comparisons across fields raise a fairness question."
%   "expand the literature beyond the current five references"
%   "use graphs to help readers understand the analysis ... interpret the visuals"
%   "Use first-person plural consistently to describe your work"
% Template: avoid a broad textbook opening; every objective must connect to a
% result and a conclusion; signpost the report without listing section titles.
% All references below were checked (Crossref) and are in 90_works_cited.tex.
% ============================================================

\subsection{Why Comparing GPAs Across Fields Raises a Fairness Question}
\label{sec:fairness}

\guide{Adam -- Paragraph 1 of 1.2: fields grade differently, and grades steer choices}{%
\textbf{Establish:} departments grade differently (Achen and Courant; Sabot and Wakeman-Linn), across many universities (Jeong et al.); low grades push students out of science majors (Ost; Rask) and students move towards leniently graded courses when grade information is visible (Bar et al.; read its abstract first); when Wellesley capped average grades in high-grading departments, enrolment in those departments fell (Butcher et al.). So grading practices and choices feed each other.}
\budget{140}{Adam}

\guide{Adam -- Paragraph 2 of 1.2: the fairness question}{%
\textbf{Establish (Irene):} universities and programmes rank students from different degrees on one GPA list (exchange places, grants, honours, master's admission); if a grade does not mean the same in every field, the list partly sorts students by the field they chose -- and, through \S\ref{sec:context}, by their background. One concrete example the reader knows (exchange-place allocation).\\
\textbf{Existing answers:} adjusted GPAs that model each course's grading (Young; Johnson); comparisons of GPA measures against later outcomes (Bailey et al. -- read the abstract first); grade distribution tables to compare grades across groups (European Commission, ECTS Users' Guide).\\
\textbf{Make clear (outline 2a):} we measure differences in \emph{observed grades}; we do not claim a field is harder or its students better.\\
\textbf{Idris's feedback:} announce that we test his adjusted mark (original mark minus the field deviation) on selection outcomes.}
\budget{160}{Adam}

